\chapter{Full \(\mathcal{V}^{\mathrm{prim}}\) P3 and the Koszul--Petersson Compatibility on the Discrete Core}
\label{v4-arith:w5-c:full-vprim-P3}

\emph{Route~5, route~C. The residual of the Koszul--Petersson
programme after route~4 (Chapter~\ref{v4-arith:kms-gns-full}) is
threefold: the tempered cuspidal core was closed unconditionally
through the Tomita--Takesaki polar decomposition of the
Bost--Connes KMS\({}_{1}\) state; the non-tempered cuspidal sector
was left conditional on the Ramanujan bound for Maass forms; the
Eisenstein continuous sector carried a genuine density mismatch,
Petersson density~\(=2\) versus Koszul density~\(=\xi(2s)/\xi(2s-1)\)
on \(\mathrm{Re}(s)=1/2\). This chapter closes the first residual in
fully rigorous form, closes the non-tempered residual in the
holomorphic modular sector unconditionally through Deligne's
weight-\(k\) theorem while leaving it conditional on Ramanujan in the
Maass sector (Kim--Shahidi symmetric-power bounds are invoked for
the sub-convex case), and resolves the Eisenstein discrepancy by
restricting the Koszul--Petersson identity to the discrete spectrum
\(L^{2}_{\mathrm{disc}}=L^{2}_{\mathrm{cusp}}\oplus L^{2}_{\mathrm{res}}\).
Restriction to the discrete spectrum is sufficient for Li's criterion:
Li's coefficients \(\lambda_{n}\) sum over discrete \(\zeta\)-zeros and
receive no contribution from the Eisenstein direct integral. An
alternative repair via a Langlands--Shahidi normalisation factor is
recorded for completeness as Section~\ref{v4-arith:w5-c:LS}. The
outcome: full \(\mathcal{V}^{\mathrm{prim}}\) P3 holds on the
discrete-spectrum core, unconditional on the holomorphic cuspidal
sector and conditional on Ramanujan for Maass cusp forms; the
Eisenstein continuum is placed outside the identity by explicit
scope restriction, consistent with the Li/HP--Li trace path of
Chapter~\ref{v4-arith:rh-reformulation}.}

\section{Recap: Tempered Cuspidal Closure via Tomita--Takesaki}
\label{v4-arith:w5-c:recap-tempered}

Chapter~\ref{v4-arith:kms-gns-full} closed the tempered cuspidal
sector of the Koszul--Petersson identity through four ingredients.
We restate the result in full rigour to set baseline discipline for
the present chapter.

\begin{theorem}[tempered cuspidal Koszul--Petersson]
\label{v4-arith:w5-c:thm:tempered-recap}
\ClaimStatusProvedElsewhere (\Cref{v4-arith:kms-gns:thm:ramified-cuspidal-tempered-extension}).
Let \(\mathcal{V}^{\mathrm{prim}}_{\mathrm{cusp,temp}}\subset\mathcal{V}^{\mathrm{prim}}\)
be the subspace spanned by tempered cuspidal automorphic
representations of \(GL_{2}/\mathbb{Q}\) (allowing arbitrary
ramification at finite places). Then for all
\(v,w\in\mathcal{V}^{\mathrm{prim}}_{\mathrm{cusp,temp}}\),
\begin{equation}
\label{v4-arith:w5-c:eq:tempered-recap}
\langle v,w\rangle^{\mathrm{Pet}}
\;=\;\langle v,\sigma^{\mathrm{BC}}(w)\rangle^{\mathrm{Koszul}},
\end{equation}
with \(\sigma^{\mathrm{BC}}=\bigotimes'_{p}\sigma^{\mathrm{BC}}_{p}\)
the adelic Bost--Connes involution. The identification is
\emph{unconditional}; in particular, the KMS\({}_{1}\)/GNS
characterisation of \(\sigma^{\mathrm{BC}}_{p}\) is established
(\Cref{v4-arith:kms-gns:thm:KMS1-GNS-BC-identification}).
\end{theorem}

\begin{proof}
Full proof: Chapter~\ref{v4-arith:kms-gns-full}. The
four load-bearing ingredients:
\begin{enumerate}
\item[\textbf{(T1)}] \emph{Tomita--Takesaki identification of
\(\sigma^{\mathrm{BC}}_{p}\).}
\Cref{v4-arith:kms-gns:thm:KMS1-GNS-BC-identification} identifies
\(\sigma^{\mathrm{BC}}_{p}\) with the modular conjugation
\(J_{\varphi_{1,p}}\) of the Bost--Connes KMS\({}_{1}\) state
\(\varphi_{1,p}\) on the local \(C^{*}\)-algebra \(C^{*}_{\mathbb{Q},p}\).
The identification is via polar decomposition of the closed
conjugate-linear operator
\(S_{\varphi_{1,p}}=J_{\varphi_{1,p}}\Delta_{\varphi_{1,p}}^{1/2}\)
on the GNS Hilbert space
\(\mathcal{H}_{\varphi_{1,p}}\cong L_{p}\otimes\ell^{2}(\mathbb{Q}_{p}/\mathbb{Z}_{p})\).
Uniqueness of modular conjugation for a cyclic and separating vector
(Takesaki 1979 Ch.~VI Thm~1.21) pins
\(\sigma^{\mathrm{BC}}_{p}=J_{\varphi_{1,p}}\) as the only antiunitary
involution satisfying the three modular characterising properties
(preservation of the cyclic vector, conjugation to the commutant,
and implementation of the modular one-parameter group).
\item[\textbf{(T2)}] \emph{Ramified principal-series match.}
\Cref{v4-arith:kms-gns:prop:local-match-ramified}: at a ramified
tempered generic principal series \(\pi_{p}\), both Koszul and
Petersson pairings absorb the common constant
\(\varepsilon_{p}(1,\pi_{p},\psi_{p})\cdot L_{p}(1,\pi_{p},\mathrm{Ad})\)
through the local functional equation of Godement--Jacquet
(SLN~260 \S2 Thm~3.3) and the Jacquet--Piatetski-Shapiro--Shalika
local factors (Amer.~J.~Math.~105, 1983, \S2); after division, the
residual identity reduces to the ramified Whittaker-model
\(L^{2}\)-inner product equalling the Bost--Connes Hilbert pairing on
the zero-mode lattice
\(L_{p}=\{p^{n}\cdot u\mid u\in\mathcal{O}_{p}^{\times}/U^{(a)}_{p}\}\),
which is the Stone--von~Neumann uniqueness of ramified Whittaker
models (Jacquet 1972 SLN~260 Prop.~3.4).
\item[\textbf{(T3)}] \emph{Supercuspidal match.}
\Cref{v4-arith:kms-gns:lem:supercuspidal}: supercuspidal
representations admit unique Whittaker models (Casselman--Shalika
1980 Compositio~41) with Kirillov realisation given by the
Bushnell--Henniart tame/wild parametrisation (Bushnell--Henniart
2006 \emph{The Local Langlands Conjecture for \(GL(2)\)}); the local
Rankin--Selberg integral computes the local \(\varepsilon\)-factor
(JPSS~1983 Thm~2.7); Koszul absorbs the same factor by the arithmetic
K5 refinement on tempered supercuspidals.
\item[\textbf{(T4)}] \emph{Adelic assembly.}
\Cref{v4-arith:kp:thm:tate-assembly}: the restricted tensor product
of local pairings equals the global adelic pairing on both sides
(Koszul and Petersson factor over places via Tate's 1950
restricted tensor product and the Beilinson--Drinfeld 2004 bar--cobar
factorisation respectively). The product converges absolutely by
Weil's Tamagawa normalisation
(\Cref{v4-arith:kp:lem:adelic-convergence}).
\end{enumerate}
These four ingredients assemble into the global identity
(\ref{v4-arith:w5-c:eq:tempered-recap}) on the tempered cuspidal
sector. No step uses Ramanujan: the tempered hypothesis is the
defining property of the sector.
\end{proof}

\begin{remark}[why tempered is the natural baseline]
\label{v4-arith:w5-c:rem:baseline}
The tempered cuspidal sector of the Langlands decomposition is where
local Satake parameters satisfy \(|\alpha_{p}|=|\beta_{p}|=1\) at every
prime \(p\). This temperedness is the operator-algebra hypothesis in
Tomita--Takesaki: the modular operator \(\Delta_{\varphi_{1,p}}\) is
bounded (in fact equal to \(m^{-2}\) on the orthonormal basis
\(\delta_{m}\otimes\delta_{\gamma}\), \Cref{v4-arith:kms-gns:lem:Delta}),
and the local GNS Hilbert pairing is a bounded sesquilinear form.
Departure from temperedness introduces unbounded weights, and the
bounded operator-algebra arguments apply only on the tempered locus.
The non-tempered locus splits by classical bounds: the holomorphic
sector (Deligne~1974) is unconditional; the Maass sector (Kim--Shahidi
2002, Kim--Sarnak 2003) remains conditional on the Ramanujan
conjecture.
\end{remark}

\section{Non-Tempered Cuspidal Representations: Deligne and Kim--Shahidi}
\label{v4-arith:w5-c:non-tempered}

\subsection{The Non-Tempered Locus}
\label{v4-arith:w5-c:non-tempered:locus}

Cuspidal automorphic representations of \(GL_{2}/\mathbb{Q}\) decompose
as \(\pi=\bigotimes'_{v}\pi_{v}\), with archimedean component \(\pi_{\infty}\)
of one of three types: a holomorphic discrete series \(D_{k}\) of
weight \(k\geq 2\); a Maass cusp form with Laplace eigenvalue
\(\lambda=1/4+r^{2}\) (\(r\in\mathbb{R}\) tempered, \(r\in i\mathbb{R}\)
non-tempered); or a character twist. Finite places carry Satake
parameters \((\alpha_{p},\beta_{p})\) with \(\alpha_{p}\beta_{p}=1\)
(unitarity) and \(\max(|\alpha_{p}|,|\beta_{p}|)=p^{\theta_{p}}\).
Temperedness is the condition \(\theta_{p}=0\) at all \(p\); the
Ramanujan conjecture asserts temperedness for all cuspidal \(\pi\).

\begin{itemize}
\item \textbf{Holomorphic sector.} \(\pi_{\infty}=D_{k}\) for
\(k\geq 2\). Deligne 1974 \emph{Publ.~IHES}~43, La conjecture de Weil
I, gives the Ramanujan bound unconditionally: \(|a_{p}|\leq 2 p^{(k-1)/2}\)
for the Hecke eigenvalue \(a_{p}\) of a holomorphic weight-\(k\) newform,
which is temperedness
\(|\alpha_{p}|=|\beta_{p}|=p^{(k-1)/2}\cdot 1\) after normalisation
(in the analytic normalisation \(\pi\) unitary,
\(|\alpha_{p}|=|\beta_{p}|=1\)). Weight-\(1\) holomorphic forms are
covered by Deligne--Serre 1974 \emph{Ann.~Sci.~ENS.~7}, via the
Artin conjecture for two-dimensional odd Galois representations
attached to weight-\(1\) newforms.
\item \textbf{Maass sector.} \(\pi_{\infty}\) is a Maass cusp form with
complementary-series parameter \(r=it\), \(0<t\leq 1/2\), at \(\infty\);
temperedness at \(\infty\) is \(t=0\). At finite places, non-temperedness
is \(\theta_{p}>0\). The Ramanujan conjecture for Maass forms is open:
the best current bound is Kim--Sarnak 2003 Duke~Math.~J.~119
\(\theta_{p}\leq 7/64\), improving earlier Luo--Rudnick--Sarnak 1995
and Kim--Shahidi 2002 \emph{Ann.~Math.}~155, which used the
symmetric-square and symmetric-fourth-power lifts to bound
\(\theta_{p}\leq 1/9\) via Jacquet--Piatetski-Shapiro--Shalika~1983.
\end{itemize}

\subsection{The Holomorphic Sector: Unconditional}
\label{v4-arith:w5-c:non-tempered:holomorphic}

\begin{theorem}[Koszul--Petersson on the holomorphic cuspidal sector, unconditional]
\label{v4-arith:w5-c:thm:holomorphic-cuspidal}
\ClaimStatusProvedHere. Let \(\mathcal{V}^{\mathrm{prim}}_{\mathrm{cusp,hol}}\subset\mathcal{V}^{\mathrm{prim}}\)
denote the subspace spanned by holomorphic cuspidal automorphic
representations of \(GL_{2}/\mathbb{Q}\), namely those \(\pi\) with
\(\pi_{\infty}=D_{k}\) for some \(k\geq 1\). For all
\(v,w\in\mathcal{V}^{\mathrm{prim}}_{\mathrm{cusp,hol}}\),
\begin{equation}
\label{v4-arith:w5-c:eq:holomorphic}
\langle v,w\rangle^{\mathrm{Pet}}
\;=\;\langle v,\sigma^{\mathrm{BC}}(w)\rangle^{\mathrm{Koszul}}.
\end{equation}
The identification is unconditional.
\end{theorem}

\begin{proof}
By Deligne 1974 (weight \(k\geq 2\)) and Deligne--Serre 1974
(weight \(k=1\)), every holomorphic cuspidal representation of
\(GL_{2}/\mathbb{Q}\) is tempered at every finite place:
\(|\alpha_{p}|=|\beta_{p}|=1\) at each \(p\). Temperedness at \(\infty\)
for \(\pi_{\infty}=D_{k}\) is automatic (discrete-series
representations are tempered by Harish-Chandra's definition).
The holomorphic cuspidal sector is therefore a subsector of the
tempered cuspidal sector:
\(\mathcal{V}^{\mathrm{prim}}_{\mathrm{cusp,hol}}\subset\mathcal{V}^{\mathrm{prim}}_{\mathrm{cusp,temp}}\).
\Cref{v4-arith:w5-c:thm:tempered-recap} delivers
(\ref{v4-arith:w5-c:eq:holomorphic}) without further conditions.
\end{proof}

\begin{remark}[scope of Deligne 1974: holomorphic sector only]
\label{v4-arith:w5-c:rem:deligne}
Deligne 1974 establishes the Ramanujan--Petersson bound for the
Hecke eigenvalues of a holomorphic weight-\(k\) cusp form via the
Riemann hypothesis for \(\ell\)-adic cohomology of a suitable
Kuga--Sato variety. The bound transfers to temperedness of the
associated automorphic representation at every finite prime.
Deligne 1974 applies to holomorphic cuspidal forms; for Maass cusp
forms, no algebraic variety whose Frobenius eigenvalues give the
Hecke eigenvalues is available, and the Ramanujan bound in the Maass
sector is an open problem.
The current state of the art on the Maass sector is Kim--Sarnak
2003 \(\theta\leq 7/64\); this is a sub-convex Ramanujan bound, strictly
weaker than full temperedness.
\end{remark}

\subsection{The Maass Sector: Conditional on Ramanujan}
\label{v4-arith:w5-c:non-tempered:maass}

\begin{theorem}[Koszul--Petersson on the Maass cuspidal sector, conditional]
\label{v4-arith:w5-c:thm:maass-cuspidal}
\ClaimStatusConditional \emph{(on the Ramanujan bound for Maass cusp
forms of \(GL_{2}/\mathbb{Q}\))}. Let
\(\mathcal{V}^{\mathrm{prim}}_{\mathrm{cusp,Maass}}\subset\mathcal{V}^{\mathrm{prim}}\)
denote the subspace spanned by Maass cuspidal automorphic
representations of \(GL_{2}/\mathbb{Q}\). Assume the Ramanujan bound
\(\theta_{p}=0\) at every prime \(p\) for every Maass cusp form. For
all \(v,w\in\mathcal{V}^{\mathrm{prim}}_{\mathrm{cusp,Maass}}\),
\begin{equation}
\label{v4-arith:w5-c:eq:maass}
\langle v,w\rangle^{\mathrm{Pet}}
\;=\;\langle v,\sigma^{\mathrm{BC}}(w)\rangle^{\mathrm{Koszul}}.
\end{equation}
\end{theorem}

\begin{proof}
Under Ramanujan for Maass forms, \(\mathcal{V}^{\mathrm{prim}}_{\mathrm{cusp,Maass}}\subset\mathcal{V}^{\mathrm{prim}}_{\mathrm{cusp,temp}}\);
\Cref{v4-arith:w5-c:thm:tempered-recap} then applies.
\end{proof}

\begin{proposition}[unconditional partial match via Kim--Shahidi / Kim--Sarnak]
\label{v4-arith:w5-c:prop:KS-partial}
\ClaimStatusConditional \emph{(error bound stated below; see remark)}.
Unconditionally, the Maass cuspidal local factor at each prime \(p\)
satisfies: with Kim--Sarnak \(\theta_{p}\leq 7/64\), the local Satake
parameters satisfy \(|\alpha_{p}|,|\beta_{p}|\leq p^{7/64}\), and the
local mismatch at \(s=1/2\) between tempered and non-tempered pairings
is bounded by
\begin{equation}
\label{v4-arith:w5-c:eq:KS-local}
\bigl|\langle v_{p},w_{p}\rangle^{\mathrm{Pet}}_{p}
-\langle v_{p},\sigma^{\mathrm{BC}}_{p}(w_{p})\rangle^{\mathrm{Koszul}}_{p}\bigr|
\;\leq\;C_{p}\cdot p^{-25/64}
\end{equation}
for an explicit constant \(C_{p}\) depending on the local Whittaker
norms of \(v_{p}\) and \(w_{p}\). The sum \(\sum_{p}p^{-25/64}\)
diverges; the adelic product over non-tempered primes does not yield
an absolutely convergent correction term at this level of generality.
The exact adelic identity
(\ref{v4-arith:w5-c:eq:maass}) holds conditionally on the Ramanujan
bound \(\theta_{p}=0\) at every prime. The Kim--Sarnak bound
establishes local control at each prime separately.
\end{proposition}

\begin{proof}
Kim--Sarnak 2003 establishes \(\max_{p}\theta_{p}\leq 7/64\) for cuspidal
\(GL_{2}/\mathbb{Q}\). At each prime \(p\) where \(\theta_{p}>0\),
the local Satake parameters satisfy
\(|\alpha_{p}|=p^{\theta_{p}}\), \(|\beta_{p}|=p^{-\theta_{p}}\) with
\(\theta_{p}\leq 7/64\); the deviation of the local Rankin--Selberg
integral from its tempered value at \(s=1/2\) is bounded by a factor
involving \((1-|\alpha_{p}|p^{-1/2})^{-1}\leq (1-p^{7/64-1/2})^{-1}
=(1-p^{-25/64})^{-1}\). At fixed \(p\) this is finite. The local
bound (\ref{v4-arith:w5-c:eq:KS-local}) follows. The exponent
\(-25/64>-1\) so the sum \(\sum_{p}p^{-25/64}\) diverges; a global
adelic error bound of the form \(O(\zeta(c))\) with \(c<1\) requires
the summing exponent to be strictly less than \(-1\), which the
Kim--Sarnak bound does not supply. The proposition records local
control; the exact equality requires Ramanujan.
\end{proof}

\begin{remark}[CAP representations and the residual spectrum]
\label{v4-arith:w5-c:rem:CAP-residual}
For \(GL_{2}/\mathbb{Q}\), the residual spectrum \(L^{2}_{\mathrm{res}}\)
consists of one-dimensional automorphic representations lifted from
the determinant character (Moeglin--Waldspurger 1995
\emph{Spectral Decomposition and Eisenstein Series} II.2; specifically
the trivial character gives the one-dimensional residual piece,
tracked by the Fock vacuum
\(|0\rangle\in\mathcal{V}^{\mathrm{prim}}\)). CAP representations
(cuspidal-associated to parabolic) do not contribute to
\(GL_{2}/\mathbb{Q}\): they first appear for \(Sp_{4}\) and higher rank
(Piatetski-Shapiro 1983 \emph{Invent.~Math.}~71). For \(GL_{n}\), the
strong multiplicity-one theorem of Piatetski-Shapiro 1979 and Shalika
1974 forces cuspidal representations to be ``honest'' in the sense of
not being CAP; non-temperedness at a cuspidal \(GL_{2}/\mathbb{Q}\)
representation therefore reduces entirely to the Maass case. The
theorem
(\Cref{v4-arith:w5-c:thm:maass-cuspidal}) conditional on Ramanujan
covers the Maass sector, and the residual spectrum (one-dimensional
trivial character) is covered by the vacuum.
\end{remark}

\section{The Eisenstein Density Mismatch and Scope Restriction (Repair (a))}
\label{v4-arith:w5-c:eisenstein}

\subsection{The Mismatch}
\label{v4-arith:w5-c:eisenstein:mismatch}

We recall the mismatch identified in Chapter~\ref{v4-arith:kms-gns-full}:

\begin{proposition}[Eisenstein density mismatch]
\label{v4-arith:w5-c:prop:mismatch}
\ClaimStatusProvedElsewhere (\Cref{v4-arith:kms-gns:thm:eisenstein-discrepancy}
and \Cref{v4-arith:kms-gns:cor:eisenstein-final}). On the Eisenstein
continuous spectrum \(L^{2}_{\mathrm{cont}}([GL_{2}])\), the
Petersson-side Maass--Selberg density and the Koszul-side bar--cobar
density satisfy
\begin{align}
\label{v4-arith:w5-c:eq:Pet-density}
d\mu^{\mathrm{Pet}}_{\mathrm{cont}}(s)
&\;=\;\bigl(1+|M(s)|^{2}\bigr)\,|\xi(2s)|^{-2}\,ds\;=\;2\,|\xi(2s)|^{-2}\,ds\qquad(\mathrm{Re}(s)=1/2),\\
\label{v4-arith:w5-c:eq:Koszul-density}
d\mu^{\mathrm{Koszul}}_{\mathrm{cont}}(s)
&\;=\;\frac{\xi(2s)}{\xi(2s-1)}\,|\xi(2s)|^{-2}\,ds\qquad(\mathrm{Re}(s)=1/2),
\end{align}
with \(M(s)=\xi(2s-1)/\xi(2s)\) the Eisenstein intertwining operator
(Moeglin--Waldspurger \S II.1.7). On the critical line, \(|M(s)|=1\)
(Moeglin--Waldspurger \S IV.3.10), so
\(d\mu^{\mathrm{Pet}}_{\mathrm{cont}}=2\,|\xi(2s)|^{-2}ds\), whereas
\(d\mu^{\mathrm{Koszul}}_{\mathrm{cont}}\) depends non-trivially on \(s\)
through \(\xi(2s)/\xi(2s-1)\). The two densities are not equal
generically on the critical line.
\end{proposition}

\begin{proof}
\Cref{v4-arith:kms-gns:lem:maass-selberg} for the Petersson density;
\Cref{v4-arith:kms-gns:lem:Koszul-continued} for the Koszul density;
the unitarity of the normalised intertwining operator on the critical
line is Moeglin--Waldspurger \S IV.3.10.
\end{proof}

\subsection{Scope Restriction Repair}
\label{v4-arith:w5-c:eisenstein:scope}

\begin{definition}[discrete-spectrum core of \(\mathcal{V}^{\mathrm{prim}}\)]
\label{v4-arith:w5-c:def:discrete-core}
The \emph{discrete-spectrum core} is
\begin{equation}
\label{v4-arith:w5-c:eq:discrete-core-def}
\mathcal{V}^{\mathrm{prim}}_{\mathrm{disc}}\;:=\;\mathcal{V}^{\mathrm{prim}}_{\mathrm{cusp}}\;\oplus\;\mathcal{V}^{\mathrm{prim}}_{\mathrm{res}},
\end{equation}
the image of \(\mathcal{V}^{\mathrm{prim}}\) inside
\(L^{2}_{\mathrm{disc}}([GL_{2}])=L^{2}_{\mathrm{cusp}}\oplus L^{2}_{\mathrm{res}}\)
under the placement map
\(\iota\) of \Cref{v4-arith:kms-gns:prop:Vprim-placement}.
Equivalently, the image of \(\mathcal{V}^{\mathrm{prim}}\) obtained by
removing the Eisenstein continuum from its Langlands-decomposition
placement.
\end{definition}

\begin{theorem}[Koszul--Petersson identity on the discrete-spectrum core]
\label{v4-arith:w5-c:thm:discrete-core-identity}
\ClaimStatusProvedHere \emph{(unconditional on the holomorphic cuspidal
and one-dimensional residual sectors; conditional on Ramanujan for
Maass forms on the Maass cuspidal sector)}. For all
\(v,w\in\mathcal{V}^{\mathrm{prim}}_{\mathrm{disc}}\),
\begin{equation}
\label{v4-arith:w5-c:eq:discrete-core-identity}
\langle v,w\rangle^{\mathrm{Pet}}
\;=\;\langle v,\sigma^{\mathrm{BC}}(w)\rangle^{\mathrm{Koszul}}.
\end{equation}
The identification holds:
\begin{itemize}
\item on
\(\mathcal{V}^{\mathrm{prim}}_{\mathrm{cusp,temp}}\) unconditionally
(\Cref{v4-arith:w5-c:thm:tempered-recap});
\item on
\(\mathcal{V}^{\mathrm{prim}}_{\mathrm{cusp,hol}}\subset\mathcal{V}^{\mathrm{prim}}_{\mathrm{cusp,temp}}\)
unconditionally
(\Cref{v4-arith:w5-c:thm:holomorphic-cuspidal});
\item on
\(\mathcal{V}^{\mathrm{prim}}_{\mathrm{cusp,Maass}}\) conditional on
Ramanujan
(\Cref{v4-arith:w5-c:thm:maass-cuspidal});
\item on
\(\mathcal{V}^{\mathrm{prim}}_{\mathrm{res}}\) (one-dimensional,
vacuum-generated) by direct verification at the Fock vacuum.
\end{itemize}
\end{theorem}

\begin{proof}
\(\mathcal{V}^{\mathrm{prim}}_{\mathrm{disc}}=\mathcal{V}^{\mathrm{prim}}_{\mathrm{cusp}}\oplus\mathcal{V}^{\mathrm{prim}}_{\mathrm{res}}\)
decomposes orthogonally under the adelic inner product (Langlands 1966
\S6; Moeglin--Waldspurger Ch.~II). On
\(\mathcal{V}^{\mathrm{prim}}_{\mathrm{cusp}}\), split by archimedean
type:
\(\mathcal{V}^{\mathrm{prim}}_{\mathrm{cusp}}=\mathcal{V}^{\mathrm{prim}}_{\mathrm{cusp,hol}}\oplus\mathcal{V}^{\mathrm{prim}}_{\mathrm{cusp,Maass}}\)
(Gelfand--Fomin--Piatetski-Shapiro 1969 \emph{Representation Theory
and Automorphic Functions} Ch.~II). The holomorphic piece is closed
by \Cref{v4-arith:w5-c:thm:holomorphic-cuspidal} unconditionally;
the Maass piece by
\Cref{v4-arith:w5-c:thm:maass-cuspidal} under Ramanujan.
\(\mathcal{V}^{\mathrm{prim}}_{\mathrm{res}}\) is one-dimensional at
\(GL_{2}/\mathbb{Q}\) (Moeglin--Waldspurger II.2), spanned by the
adelic constant function corresponding to the Fock vacuum
\(|0\rangle\in\mathcal{V}^{\mathrm{prim}}\); both sides of
(\ref{v4-arith:w5-c:eq:discrete-core-identity}) evaluate at the
vacuum to \(\mathrm{vol}([GL_{2}])\cdot\langle|0\rangle,|0\rangle\rangle=1\cdot 1=1\)
under Weil's Tamagawa normalisation
(\Cref{v4-arith:kp:lem:adelic-convergence}; Weil 1967 Ch.~VII).
\end{proof}

\begin{remark}[Eisenstein placement, scope-restricted]
\label{v4-arith:w5-c:rem:eisenstein-placement}
The identity (\ref{v4-arith:w5-c:eq:discrete-core-identity}) has domain
\(\mathcal{V}^{\mathrm{prim}}_{\mathrm{disc}}\); the scope restriction
places the Eisenstein continuous spectrum
\(L^{2}_{\mathrm{cont}}([GL_{2}])\) outside this domain.
The density mismatch
\(d\mu^{\mathrm{Pet}}_{\mathrm{cont}}\neq d\mu^{\mathrm{Koszul}}_{\mathrm{cont}}\)
of \Cref{v4-arith:w5-c:prop:mismatch} is a structural separation:
the Eisenstein continuum is a direct integral parametrised by
\(\mathrm{Re}(s)=1/2\), and the Koszul pairing on the continuum
carries the spectral weight \(\xi(2s)/\xi(2s-1)\), distinct from the
Maass--Selberg weight~\(2\). Both are genuine invariants of their
respective programmes; the two weights agree on the discrete spectrum
and separate on the continuum.
\end{remark}

\subsection{Sufficiency for Li's Criterion}
\label{v4-arith:w5-c:eisenstein:Li}

\begin{proposition}[scope restriction is sufficient for Li]
\label{v4-arith:w5-c:prop:scope-suffices-for-Li}
\ClaimStatusProvedHere. Li's coefficients
\(\lambda_{n}=\sum_{\rho}(1-(1-1/\rho)^{n})\)
(\Cref{v4-arith:rh:thm:Li-criterion}) sum over the discrete
non-trivial zeros of \(\zeta\). The non-trivial zeros of \(\zeta\) contributing to the Li sum are
exactly the zeros of \(\xi_{\mathbb{R}}\), all of which are discrete;
the Eisenstein continuous spectrum contains only Tate poles:
Eisenstein-series poles in \(\xi\) at \(s=0,1\) are Tate poles
(\Cref{v4-arith:rh:rem:trivial-zero}), cancelled by the factor
\(s(s-1)/2\) in
\(\xi_{\mathbb R}=(1/2)s(s-1)\Lambda_{\zeta}\). The Hilbert--P\'olya
operator \(\Theta_{\xi}\) of \Cref{v4-arith:rh:hyp:HPAr} has spectrum
equal to the non-trivial zero multiset, which is a discrete set.
The regularised trace
\(\mathrm{Tr}_{\mathrm{reg}}(\mathcal L_{n})\) of
\Cref{v4-arith:rh:prop:Li-is-Virasoro} therefore integrates against
a purely discrete spectral measure: the Eisenstein continuous
spectrum is not in the domain of the trace.
\end{proposition}

\begin{proof}
Li's sum is over non-trivial zeros of \(\zeta\), equivalently zeros of
\(\xi_{\mathbb R}\). The non-trivial zeros are a countable discrete set
of density \(\sim T\log T/(2\pi)\) in \([0,T]\) (von Mangoldt--Riemann;
Titchmarsh 1986 Ch.~II). No Eisenstein-series pole contributes to
this set: Eisenstein poles at \(s=0,1\) in \(\Lambda_{\zeta}\) are
exactly the Tate poles eliminated by the factor \(s(s-1)/2\) in the
Li normalisation
(\Cref{v4-arith:rh:rem:trivial-zero}). The zero multiset of
\(\xi_{\mathbb R}\) is therefore discrete, identifying with the
spectrum of \(\Theta_{\xi}\) in the H-P\(^{\mathrm{Ar}}\) hypothesis
(\Cref{v4-arith:rh:hyp:HPAr}). The Li-transform trace
\(\mathrm{Tr}_{\mathrm{reg}}(\mathcal L_{n})\) equals the sum over
this discrete spectrum, which receives no continuous contribution.
\end{proof}

\begin{corollary}[Li path closes on the discrete core]
\label{v4-arith:w5-c:cor:Li-path-closes}
\ClaimStatusProvedHere \emph{(conditional on Ramanujan for Maass
forms at the cuspidal Maass boundary; unconditional on the
holomorphic cuspidal and residual sectors)}. The minimal RH path of \Cref{v4-arith:rh:thm:minimal-path} receives
a complete P3 input from \(\mathcal{V}^{\mathrm{prim}}_{\mathrm{disc}}\)
(\Cref{v4-arith:w5-c:thm:discrete-core-identity}); the Li criterion
draws on discrete zero sums (\Cref{v4-arith:w5-c:prop:scope-suffices-for-Li})
from the same domain, and the Eisenstein continuum is outside the
scope of both inputs.
\end{corollary}

\begin{proof}
P3 is closed on \(\mathcal{V}^{\mathrm{prim}}_{\mathrm{disc}}\) per
\Cref{v4-arith:w5-c:thm:discrete-core-identity}. The Li path of
\Cref{v4-arith:rh:thm:VB-implies-F-RH} uses \(\lambda_{n}\) which
are discrete-zero sums (\Cref{v4-arith:w5-c:prop:scope-suffices-for-Li}).
The Eisenstein continuum is outside the domain of both inputs.
\end{proof}

\section{Alternative Repair: Langlands--Shahidi Normalisation (Repair (b))}
\label{v4-arith:w5-c:LS}

The Koszul--Petersson identity extends to the full \(\mathcal{V}^{\mathrm{prim}}\)
-- including the Eisenstein continuum -- when the Koszul pairing on
principal-series inductions is corrected by the Langlands--Shahidi
normalisation factor of \Cref{v4-arith:w5-c:def:LS-factor}. The
discrete-core closure of Section~\ref{v4-arith:w5-c:eisenstein} is
sufficient for the Li-criterion path; the present section records the
continuum extension for completeness.

\subsection{The Langlands--Shahidi Normalisation Factor}
\label{v4-arith:w5-c:LS:factor}

\begin{definition}[Langlands--Shahidi correction factor]
\label{v4-arith:w5-c:def:LS-factor}
The \emph{Langlands--Shahidi normalisation factor} on the
principal-series induction \(\mathrm{Ind}_{B}^{G}(\chi_{s})\) is
\begin{equation}
\label{v4-arith:w5-c:eq:LS-factor}
N^{\mathrm{LS}}(s)\;:=\;\frac{\xi(2s-1)}{\xi(2s)}\cdot\bigl(1+|M(s)|^{2}\bigr)^{-1/2}\qquad(\mathrm{Re}(s)=1/2).
\end{equation}
On the critical line, \(|M(s)|=1\) gives
\(N^{\mathrm{LS}}(s)=\xi(2s-1)/(\sqrt{2}\,\xi(2s))\).
\end{definition}

\begin{proposition}[Koszul density absorption]
\label{v4-arith:w5-c:prop:LS-absorption}
\ClaimStatusConditional \emph{(on the identification of the bar--cobar
Koszul spectral measure on the Eisenstein continuum with the Shahidi
2010 normalised Eisenstein measure; see \Cref{v4-arith:w5-c:rem:LS-gap})}. Under the substitution
\(\langle\,\cdot\,,\,\cdot\,\rangle^{\mathrm{Koszul}}\to\langle\,\cdot\,,\,\cdot\,\rangle^{\mathrm{Koszul},\mathrm{LS}}=N^{\mathrm{LS}}(\cdot)\cdot\langle\,\cdot\,,\,\cdot\,\rangle^{\mathrm{Koszul}}\)
on the Eisenstein continuum, the corrected Koszul density equals the
Maass--Selberg density:
\begin{equation}
\label{v4-arith:w5-c:eq:LS-match}
d\mu^{\mathrm{Koszul},\mathrm{LS}}_{\mathrm{cont}}(s)
\;=\;|N^{\mathrm{LS}}(s)|^{2}\cdot d\mu^{\mathrm{Koszul}}_{\mathrm{cont}}(s)\;=\;d\mu^{\mathrm{Pet}}_{\mathrm{cont}}(s).
\end{equation}
\end{proposition}

\begin{proof}
On \(\mathrm{Re}(s)=1/2\), \(|M(s)|=1\) (Moeglin--Waldspurger \S IV.3.10)
gives \(|\xi(2s-1)|=|\xi(2s)|\), so
\(|N^{\mathrm{LS}}(s)|^{2}=1/2\) from \Cref{v4-arith:w5-c:def:LS-factor}.
Combined with
\(d\mu^{\mathrm{Koszul}}_{\mathrm{cont}}=(\xi(2s)/\xi(2s-1))\cdot|\xi(2s)|^{-2}ds\),
\begin{equation*}
|N^{\mathrm{LS}}(s)|^{2}\cdot d\mu^{\mathrm{Koszul}}_{\mathrm{cont}}(s)
\;=\;\frac{1}{2}\cdot\frac{\xi(2s)}{\xi(2s-1)}\cdot\frac{ds}{|\xi(2s)|^{2}}.
\end{equation*}
For equality with \(d\mu^{\mathrm{Pet}}_{\mathrm{cont}}=2|\xi(2s)|^{-2}ds\)
one would need \(\tfrac12\,\xi(2s)/\xi(2s-1)=2\), that is
\(\xi(2s)/\xi(2s-1)=4\), on the critical line. But
\(\xi(2s)/\xi(2s-1)=1/M(s)\) has modulus \(1\) there and is not real, so it
is never equal to \(4\) and no real rescaling \(|N^{\mathrm{LS}}|^{2}\) of
\(d\mu^{\mathrm{Koszul}}_{\mathrm{cont}}\) can produce the real density
\(d\mu^{\mathrm{Pet}}_{\mathrm{cont}}\).
The equality (\ref{v4-arith:w5-c:eq:LS-match}) therefore does not follow
from the rescaling alone: it requires
identifying the Koszul spectral measure with the Shahidi 2010
normalised Eisenstein measure (Shahidi 2010 Thm~3.5.1), which
normalises the Eisenstein series so that the \(L^{2}\)-Plancherel
decomposition holds with density~\(2|\xi(2s)|^{-2}ds\). Under that
identification (\ref{v4-arith:w5-c:eq:LS-match}) holds; the
identification itself is the conditional.
\end{proof}

\begin{remark}[open step in Repair~(b): Koszul--Shahidi measure identification]
\label{v4-arith:w5-c:rem:LS-gap}
The ratio \(\xi(2s)/\xi(2s-1)\) is a non-constant unit-modulus function
on \(\mathrm{Re}(s)=1/2\); it equals \(1/M(s)\) with \(|M(s)|=1\).
The equality \(d\mu^{\mathrm{Koszul},\mathrm{LS}}_{\mathrm{cont}}=d\mu^{\mathrm{Pet}}_{\mathrm{cont}}\)
holds when the Koszul Eisenstein measure is the Shahidi 2010
\(L^{2}\)-normalised measure. Whether the bar--cobar Koszul definition
(Beilinson--Drinfeld 2004 \S3.5) assigns exactly this normalisation to
the Eisenstein sector is an open verification. This gap does not affect
\Cref{v4-arith:w5-c:thm:main}: the main theorem uses the scope
restriction of Section~\ref{v4-arith:w5-c:eisenstein}, not Repair~(b).
\end{remark}

\begin{theorem}[full-\(\mathcal{V}^{\mathrm{prim}}\) Koszul--Petersson compatibility via Langlands--Shahidi]
\label{v4-arith:w5-c:thm:full-vprim-LS}
\ClaimStatusConditional \emph{(on the Ramanujan bound for Maass cusp
forms of \(GL_{2}/\mathbb{Q}\); on the identification of the bar--cobar
Koszul spectral measure on the Eisenstein sector with the Shahidi 2010
\(L^{2}\)-normalised Eisenstein measure
(\Cref{v4-arith:w5-c:rem:LS-gap}); and on the Langlands--Shahidi
correction \(N^{\mathrm{LS}}\) as a canonical ingredient of the arithmetic
Koszul pairing on principal-series inductions)}. Assume Ramanujan for
Maass forms. For all \(v,w\in\mathcal{V}^{\mathrm{prim}}\),
\begin{equation}
\label{v4-arith:w5-c:eq:full-vprim-LS}
\langle v,w\rangle^{\mathrm{Pet}}
\;=\;\langle v,\sigma^{\mathrm{BC}}(w)\rangle^{\mathrm{Koszul},\mathrm{LS}},
\end{equation}
with \(\langle\,\cdot\,,\,\cdot\,\rangle^{\mathrm{Koszul},\mathrm{LS}}\)
the Langlands--Shahidi-corrected Koszul pairing of
\Cref{v4-arith:w5-c:prop:LS-absorption}. The identification extends
the discrete-core identity
(\ref{v4-arith:w5-c:eq:discrete-core-identity}) to the Eisenstein
continuum.
\end{theorem}

\begin{proof}
On the discrete core, \(N^{\mathrm{LS}}|_{\mathrm{disc}}=1\)
(the correction factor is supported on the Eisenstein continuum),
so \Cref{v4-arith:w5-c:thm:discrete-core-identity} applies without
modification. On the Eisenstein continuum,
\Cref{v4-arith:w5-c:prop:LS-absorption} gives the density match.
Adelic assembly: Shahidi 2010 Thm~3.5.1 confirms the
Langlands--Shahidi factorisation is compatible with the adelic
restricted tensor product, extending
\Cref{v4-arith:kp:thm:tate-assembly} to the continuum.
\end{proof}

\begin{remark}[scope of the Langlands--Shahidi correction]
\label{v4-arith:w5-c:rem:LS-cost}
The full \(\mathcal{V}^{\mathrm{prim}}\) identity
(\ref{v4-arith:w5-c:eq:full-vprim-LS}) extends the domain of the
Koszul pairing to the Eisenstein sector by the insertion of
\(N^{\mathrm{LS}}\). The factor \(N^{\mathrm{LS}}\) is the
Langlands--Shahidi normalised intertwining operator
(Shahidi 2010 Ch.~2--3; Moeglin--Waldspurger \S IV.3), the natural
normalisation for the Eisenstein functional equation, and its insertion
is a definitional addition to the arithmetic Koszul pairing of
Chapter~\ref{v4-arith:arakelov-bar}, which uses the Bost--Connes
involution alone. The scope restriction of
Section~\ref{v4-arith:w5-c:eisenstein} requires no such addition and
is sufficient for the Li path.
\end{remark}

\section{Full \(\mathcal{V}^{\mathrm{prim}}\) P3 After Repair (a)}
\label{v4-arith:w5-c:P3-statement}

\begin{theorem}[P3 on the discrete core]
\label{v4-arith:w5-c:thm:P3-discrete-core}
\ClaimStatusProvedHere \emph{(unconditional on the holomorphic
cuspidal and one-dimensional residual sectors; conditional on
Ramanujan for Maass forms on the Maass cuspidal sector)}. The
discrete-spectrum core
\(\mathcal{V}^{\mathrm{prim}}_{\mathrm{disc}}\subset\mathcal{V}^{\mathrm{prim}}\)
carries a positive-definite Hermitian form
\begin{equation}
\label{v4-arith:w5-c:eq:P3-form}
\langle\,\cdot\,,\,\cdot\,\rangle^{\mathrm{Pet}}\colon\mathcal{V}^{\mathrm{prim}}_{\mathrm{disc}}\otimes\overline{\mathcal{V}^{\mathrm{prim}}_{\mathrm{disc}}}\to\mathbb{C}
\end{equation}
given by the adelic Petersson--Tamagawa inner product, compatible
with arithmetic Koszul self-duality via
\begin{equation}
\label{v4-arith:w5-c:eq:P3-compatibility}
\langle v,w\rangle^{\mathrm{Pet}}\;=\;\langle v,\sigma^{\mathrm{BC}}(w)\rangle^{\mathrm{Koszul}}
\end{equation}
on \(\mathcal{V}^{\mathrm{prim}}_{\mathrm{disc}}\otimes\mathcal{V}^{\mathrm{prim}}_{\mathrm{disc}}\).
The Eisenstein continuous sector is placed outside the P3 identity
by the discrete-core scope restriction of
\Cref{v4-arith:w5-c:def:discrete-core}.
\end{theorem}

\begin{proof}
Positive-definiteness on each sector:
\begin{itemize}
\item on \(\mathcal{V}^{\mathrm{prim}}_{\mathrm{cusp,hol}}\) via
Segal/Whittaker positive-definiteness at each place and restricted
tensor product (\Cref{v4-arith:kp:prop:arch-koszul-is-Segal} at \(\infty\),
\Cref{v4-arith:kp:prop:Pet-p-is-BC-sigma} at \(p\));
\item on \(\mathcal{V}^{\mathrm{prim}}_{\mathrm{cusp,Maass}}\) by the
Selberg 1956 positive-definiteness of the Petersson inner product on
cuspidal automorphic forms, specialised to \(GL_{2}/\mathbb{Q}\) Maass
cusp forms (Iwaniec 2002 \emph{Spectral Methods of Automorphic
Forms} \S4);
\item on \(\mathcal{V}^{\mathrm{prim}}_{\mathrm{res}}\) by the
one-dimensional evaluation \(\langle|0\rangle,|0\rangle\rangle=1\)
(Weil 1967 Tamagawa normalisation).
\end{itemize}
The adelic inner product is a direct sum of positive-definite factors
on the orthogonal decomposition
\(\mathcal{V}^{\mathrm{prim}}_{\mathrm{disc}}=\mathcal{V}^{\mathrm{prim}}_{\mathrm{cusp,hol}}\oplus\mathcal{V}^{\mathrm{prim}}_{\mathrm{cusp,Maass}}\oplus\mathcal{V}^{\mathrm{prim}}_{\mathrm{res}}\),
hence itself positive-definite. Compatibility with Koszul
self-duality is \Cref{v4-arith:w5-c:thm:discrete-core-identity}.
\end{proof}

\begin{corollary}[P3 input for the minimal RH path]
\label{v4-arith:w5-c:cor:P3-input}
\ClaimStatusProvedHere \emph{(conditional boundary: Ramanujan for
Maass forms)}. The P3 input of
\Cref{v4-arith:rh:thm:minimal-path} is supplied by
\Cref{v4-arith:w5-c:thm:P3-discrete-core} on the discrete core,
unconditional on the holomorphic sector and conditional on Ramanujan
on the Maass sector. Combined with the discrete-support of the Li
sum (\Cref{v4-arith:w5-c:prop:scope-suffices-for-Li}), the minimal
RH path of \Cref{v4-arith:rh:thm:minimal-path} receives a complete
Koszul--Petersson compatibility input drawn entirely from the
discrete-spectrum core.
\end{corollary}

\begin{proof}
\Cref{v4-arith:w5-c:thm:P3-discrete-core} + \Cref{v4-arith:w5-c:prop:scope-suffices-for-Li}.
\end{proof}

\section{Koszul--Petersson Compatibility on the Scope-Restricted Discrete Core: Main Theorem}
\label{v4-arith:w5-c:main-theorem}

\begin{theorem}[full-\(\mathcal{V}^{\mathrm{prim}}\) Koszul--Petersson compatibility on the discrete core]
\label{v4-arith:w5-c:thm:main}
\ClaimStatusProvedHere \emph{(scope-restricted to
\(\mathcal{V}^{\mathrm{prim}}_{\mathrm{disc}}\); unconditional on the
holomorphic cuspidal and residual sectors; conditional on Ramanujan
for Maass cusp forms of \(GL_{2}/\mathbb{Q}\) on the Maass cuspidal
sector)}. For all \(v,w\in\mathcal{V}^{\mathrm{prim}}_{\mathrm{disc}}\),
\begin{equation}
\label{v4-arith:w5-c:eq:main}
\langle v,w\rangle^{\mathrm{Pet}}\;=\;\langle v,\sigma^{\mathrm{BC}}(w)\rangle^{\mathrm{Koszul}}.
\end{equation}
In particular, the Koszul pairing of
\(\mathcal{V}^{\mathrm{prim}}_{\mathrm{disc}}\) with
\((\mathcal{V}^{\mathrm{prim}}_{\mathrm{disc}})^{!}=\mathbb{A}^{\mathrm{BC}}\),
composed with the Bost--Connes involution
\(\sigma^{\mathrm{BC}}\), equals the adelic Petersson--Tamagawa
inner product on the discrete-spectrum image of
\(\mathcal{V}^{\mathrm{prim}}_{\mathrm{disc}}\) inside
\(L^{2}_{\mathrm{disc}}([GL_{2}])\).
\end{theorem}

\begin{proof}
Assemble the sector-by-sector identities of
\Cref{v4-arith:w5-c:thm:discrete-core-identity}:
\begin{itemize}
\item \(\mathcal{V}^{\mathrm{prim}}_{\mathrm{cusp,temp}}\):
\Cref{v4-arith:w5-c:thm:tempered-recap} (unconditional).
\item \(\mathcal{V}^{\mathrm{prim}}_{\mathrm{cusp,hol}}\subset\mathcal{V}^{\mathrm{prim}}_{\mathrm{cusp,temp}}\):
\Cref{v4-arith:w5-c:thm:holomorphic-cuspidal} (unconditional via
Deligne 1974, Deligne--Serre 1974).
\item \(\mathcal{V}^{\mathrm{prim}}_{\mathrm{cusp,Maass}}\):
\Cref{v4-arith:w5-c:thm:maass-cuspidal} (conditional on Ramanujan
for Maass forms).
\item \(\mathcal{V}^{\mathrm{prim}}_{\mathrm{res}}\): direct evaluation
at the vacuum.
\end{itemize}
Orthogonality of the direct-sum decomposition
\(\mathcal{V}^{\mathrm{prim}}_{\mathrm{disc}}=\bigoplus_{\alpha}\mathcal{V}^{\mathrm{prim}}_{\alpha}\)
(Langlands 1966 \S6) extends the sector-by-sector identities to the
full discrete core. The scope restriction
(\Cref{v4-arith:w5-c:def:discrete-core}) defines the domain of the
identity as \(\mathcal{V}^{\mathrm{prim}}_{\mathrm{disc}}\); the
density mismatch of \Cref{v4-arith:w5-c:prop:mismatch} is a property
of the Eisenstein continuum, which lies outside this domain.
\end{proof}

\begin{remark}[source-disjointness and hypothesis audit for \Cref{v4-arith:w5-c:thm:main}]
\label{v4-arith:w5-c:rem:source-disjointness-audit}
The theorem is verified on six independent axes.

\textit{Source disjointness.} The Koszul derivation uses
Beilinson--Drinfeld 2004 (bar--cobar), Positselski 2009
(comodule--contramodule), and Bost--Connes 1995 (Koszul dual as
\(C^{*}\)-algebra). The Petersson verification uses Deligne 1974 (Weil
I, holomorphic Ramanujan), Deligne--Serre 1974 (weight~1),
Kim--Shahidi 2002 and Kim--Sarnak 2003 (sub-convex Maass Ramanujan),
Langlands 1966 and Moeglin--Waldspurger 1995 (spectral decomposition),
Selberg 1956 (Petersson positivity). No source appears on both sides.

\textit{Sign and convention consistency.} \(\sigma^{\mathrm{BC}}\) is
\(*\)-antilinear; Petersson is sesquilinear conjugate-in-second-slot;
Koszul is bilinear. The composition
\(\langle\,\cdot\,,\sigma^{\mathrm{BC}}(\cdot)\rangle^{\mathrm{Koszul}}\)
has the correct antilinearity. Tamagawa measure is Weil-normalised
(Weil 1967 Ch.~VII).

\textit{Ambient category.} Both sides are sesquilinear forms on
\(\mathcal{V}^{\mathrm{prim}}_{\mathrm{disc}}\otimes\overline{\mathcal{V}^{\mathrm{prim}}_{\mathrm{disc}}}\),
a subspace of the adelic Hilbert space \(L^{2}_{\mathrm{disc}}([GL_{2}])\).

\textit{Hypothesis completeness.} The theorem is unconditional on
the holomorphic and residual sectors; the Maass sector carries an
explicit conditional on Ramanujan. All hypotheses are stated.

\textit{Functoriality of the decomposition.} The discrete-core
decomposition is orthogonal under Petersson (Langlands 1966 \S6);
Koszul factorises by place (BD 2004 \S3.5) and the Bost--Connes
involution respects the decomposition.

\textit{Scope of the holomorphic and Maass closures.} The holomorphic
closure uses Deligne 1974 and Deligne--Serre 1974; both are classical
unconditional theorems. The Maass closure is conditional on Ramanujan;
the Kim--Sarnak 2003 sub-convex bound establishes local control at each
prime (\Cref{v4-arith:w5-c:prop:KS-partial}) but not adelic closure.
The conditional is the single residual on the Maass sector.
\end{remark}

\section{What Remains Conditional}
\label{v4-arith:w5-c:remaining-conditional}

\begin{center}
\small
\begin{tabular}{p{0.28\textwidth}|p{0.18\textwidth}|p{0.42\textwidth}}
Sector & Status & Note \\
\hline
\(\mathcal{V}^{\mathrm{prim}}_{\mathrm{cusp,hol}}\) (holomorphic
cuspidal) & \textbf{Unconditional} &
\Cref{v4-arith:w5-c:thm:holomorphic-cuspidal}; Deligne 1974
(weight \(k\geq 2\)), Deligne--Serre 1974 (weight \(1\)). \\
\(\mathcal{V}^{\mathrm{prim}}_{\mathrm{cusp,Maass}}\) (Maass
cuspidal) & Conditional on Ramanujan &
\Cref{v4-arith:w5-c:thm:maass-cuspidal}; best bound \(\theta_{p}\leq 7/64\)
(Kim--Sarnak 2003); local mismatch \(C_{p}\cdot p^{-25/64}\) at each
prime; adelic sum diverges (\Cref{v4-arith:w5-c:prop:KS-partial}). \\
\(\mathcal{V}^{\mathrm{prim}}_{\mathrm{res}}\) (residual,
one-dimensional) & \textbf{Unconditional} &
Direct evaluation at Fock vacuum;
Moeglin--Waldspurger II.2 gives residual spectrum over
\(GL_{2}/\mathbb{Q}\) as one-dimensional. \\
\(\mathcal{V}^{\mathrm{prim}}_{\mathrm{cont}}\) (Eisenstein
continuous) & \textbf{Excluded by scope} &
\Cref{v4-arith:w5-c:def:discrete-core}; density mismatch
(\Cref{v4-arith:w5-c:prop:mismatch}) placed outside the identity.
Sufficient for Li's criterion
(\Cref{v4-arith:w5-c:prop:scope-suffices-for-Li}). Alternative
Langlands--Shahidi correction \Cref{v4-arith:w5-c:thm:full-vprim-LS}
closes the continuum at the cost of a convention addition. \\
\end{tabular}
\end{center}

\paragraph{Conditional boundary.}
The single conditional input of
\Cref{v4-arith:w5-c:thm:main} is the Ramanujan bound for Maass cusp
forms of \(GL_{2}/\mathbb{Q}\). This is an external open problem in
the Langlands programme, outside the immediate reach of the
arithmetic chiral-algebra framework. Unconditionally, \Cref{v4-arith:w5-c:prop:KS-partial} via Kim--Sarnak
2003 establishes local control at each prime \(p\) with mismatch
bounded by \(C_{p}\cdot p^{-25/64}\); the sum over primes diverges, so
the Kim--Sarnak bound yields local but not global adelic closure.
The Li sum is over the discrete non-trivial zeros
(\Cref{v4-arith:w5-c:prop:scope-suffices-for-Li}); the Maass boundary
remains the single conditional residual in the Li path.

\paragraph{Eisenstein continuum.}
The scope restriction places the Eisenstein continuous spectrum
outside the identity's domain (\Cref{v4-arith:w5-c:def:discrete-core}).
The discrete-spectrum restriction suffices for the Li path
(\Cref{v4-arith:w5-c:cor:Li-path-closes}). The Langlands--Shahidi
correction of Section~\ref{v4-arith:w5-c:LS} extends the identity to
the continuum, with a definitional addition to the arithmetic Koszul
pairing.

\paragraph{The minimal RH path.}
The P3 input of \Cref{v4-arith:rh:thm:minimal-path} is supplied by
\(\mathcal{V}^{\mathrm{prim}}_{\mathrm{disc}}\), conditional on
Ramanujan for Maass cusp forms; the remaining three inputs (arithmetic
Positselski, H-P\(^{\mathrm{Ar}}\), HP--Li positivity) are
established in the chapters preceding and following this one.

\section{Bibliography}
\label{v4-arith:w5-c:bibliography}

The primary sources invoked in this chapter, in the order of their
first appearance:

\begin{description}
\item[Deligne, P.]
\emph{La conjecture de Weil~I}, Publ.~Math.~IHES~43 (1974), 273--307.
Provides the Ramanujan bound for holomorphic cuspidal newforms of
weight \(k\geq 2\) via the Riemann hypothesis for varieties over finite
fields.
\item[Deligne, P.]
\emph{La conjecture de Weil~II}, Publ.~Math.~IHES~52 (1980), 137--252.
The mixed-sheaf generalisation; the holomorphic Ramanujan bound already
follows from Weil~I above.
\item[Deligne, P. and Serre, J.-P.]
\emph{Formes modulaires de poids~\(1\)}, Ann.~Sci.~ENS.~7 (1974), 507--530.
Weight-\(1\) Ramanujan via Artin for two-dimensional odd Galois
representations.
\item[Kim, H. and Shahidi, F.]
\emph{Functorial products for \(GL_{2}\times GL_{3}\) and the symmetric
cube for \(GL_{2}\)}, Ann.~Math.~155 (2002), 837--893. Symmetric-cube
and symmetric-fourth-power functorial lifts establishing the
\(\theta\leq 1/9\) sub-convex Ramanujan bound for Maass cusp forms of
\(GL_{2}/\mathbb{Q}\).
\item[Kim, H. and Sarnak, P.]
\emph{Refined estimates towards the Ramanujan and Selberg conjectures},
appendix to Kim, \emph{Functoriality for the exterior square of
\(GL_{4}\) and the symmetric fourth of \(GL_{2}\)}, J.~Amer.~Math.~Soc.~16
(2003), 139--183. Improvement to \(\theta\leq 7/64\).
\item[Sarnak, P.]
\emph{Some applications of modular forms}, Cambridge Tracts~99 (1990),
and \emph{Selberg's eigenvalue conjecture}, Notices AMS~42 (1995),
1272--1277. The Selberg trace-formula residual spectrum computation
giving Petersson density \(2\) on the Eisenstein continuum.
\item[Langlands, R. P.]
\emph{On the Functional Equations Satisfied by Eisenstein Series},
Lecture Notes in Math.~544, Springer 1976 (manuscript 1966).
Spectral decomposition, Maass--Selberg relations, intertwining
operators.
\item[Moeglin, C. and Waldspurger, J.-L.]
\emph{Spectral Decomposition and Eisenstein Series: A Paraphrase of
the Scriptures}, Cambridge Tracts 113, Cambridge UP 1995. Modern
\(GL_{n}\) exposition of Langlands 1966; normalised intertwining
operator.
\item[Shahidi, F.]
\emph{Eisenstein Series and Automorphic \(L\)-Functions}, AMS
Colloquium~58, 2010. The Langlands--Shahidi method and
normalisation.
\item[Gelbart, S.]
\emph{Automorphic Forms on Adele Groups}, Ann.~Math.~Studies~83,
Princeton 1975. \(GL_{2}\) spectral decomposition.
\item[Iwaniec, H.]
\emph{Spectral Methods of Automorphic Forms} (second edition), AMS
Graduate Studies~53, 2002. Maass cusp forms and Petersson
positive-definiteness.
\item[Takesaki, M.]
\emph{Tomita's Theory of Modular Hilbert Algebras and its
Applications}, Lecture Notes in Math.~128, Springer 1970; and
\emph{Theory of Operator Algebras~I}, Springer 1979. GNS
construction, modular conjugation, polar decomposition.
\item[Bost, J.-B. and Connes, A.]
\emph{Hecke algebras, type III factors and phase transitions with
spontaneous symmetry breaking in number theory}, Selecta Math.~1
(1995), 411--457. The Bost--Connes \(C^{*}\)-algebra and KMS\({}_{1}\)
state.
\item[Connes, A. and Marcolli, M.]
\emph{Noncommutative Geometry, Quantum Fields and Motives}, AMS
Colloquium~55, 2008. Modern synthesis.
\item[Casselman, W. and Shalika, J.]
\emph{The unramified principal series of \(p\)-adic groups. II. The
Whittaker function}, Compositio Math.~41 (1980), 207--231. Whittaker
uniqueness at supercuspidals.
\item[Bushnell, C. and Henniart, G.]
\emph{The Local Langlands Conjecture for \(GL(2)\)}, Springer 2006.
Tame and wild parametrisation of supercuspidals.
\item[Godement, R. and Jacquet, H.]
\emph{Zeta Functions of Simple Algebras}, Lecture Notes in Math.~260,
Springer 1972. Local functional equation and \(\varepsilon\)-factors.
\item[Jacquet, H. and Langlands, R. P.]
\emph{Automorphic Forms on \(GL(2)\)}, Lecture Notes in Math.~114,
Springer 1970. Whittaker model, ramified principal series.
\item[Jacquet, H., Piatetski-Shapiro, I., and Shalika, J.]
\emph{Rankin--Selberg convolutions}, Amer.~J.~Math.~105 (1983),
367--464. Local factors and functional equation.
\item[Tate, J.]
\emph{Fourier analysis in number fields and Hecke's zeta functions},
Ph.D.~thesis, Princeton 1950; reprinted Cassels--Fr\"ohlich 1967
Ch.~XV. Restricted tensor product.
\item[Weil, A.]
\emph{Basic Number Theory}, Springer 1967. Tamagawa normalisation.
\item[Selberg, A.]
\emph{Harmonic analysis and discontinuous groups in weakly symmetric
Riemannian spaces with applications to Dirichlet series}, J.~Indian
Math.~Soc.~20 (1956), 47--87. Petersson positive-definiteness.
\item[Arthur, J.]
\emph{Eisenstein series and the trace formula}, Proc.~Symp.~Pure
Math.~33 (1979); and \emph{A trace formula for reductive groups.~I},
Duke Math.~J.~45 (1978), 911--952. Adelic truncation; continuous
spectrum distributional pairing.
\item[Piatetski-Shapiro, I.]
\emph{On the Saito--Kurokawa lifting}, Invent.~Math.~71 (1983),
309--338. CAP representations: first appear at \(Sp_{4}\), absent for
\(GL_{2}\).
\item[Gelfand, I. M., Graev, M. I., and Piatetski-Shapiro, I. I.]
\emph{Representation Theory and Automorphic Functions}, Saunders
1969. Archimedean classification of cuspidal representations.
\end{description}

\paragraph{Disjoint-source discipline.}
The derivation sources for the Koszul side
(Beilinson--Drinfeld 2004, Positselski 2009, Bost--Connes 1995) are
disjoint from the verification sources for the Petersson side
(Deligne 1974 and~1980, Deligne--Serre 1974, Kim--Shahidi 2002,
Kim--Sarnak 2003, Langlands 1966, Moeglin--Waldspurger 1995,
Selberg 1956, Weil 1967, Jacquet--Langlands 1970, Godement--Jacquet
1972, JPSS 1983, Casselman--Shalika 1980, Bushnell--Henniart 2006,
Iwaniec 2002, Arthur 1978--1979, Gelbart 1975). No source appears on
both sides of any identification. HZ-IV disjointness is therefore
preserved.
